\begin{lemma}
Let $n\ge1$ be an integer and let $1<p<\infty$.  Then
\begin{enumerate}
\item[(23)] $\lambda_n(n,p)=0$ and $\epsilon_n(n,p)=0$;
\item[(24)] $|\epsilon_k(n,p)|<1$ for $k=n+1,n+2,\dots$;
\item[(25)] $\dfrac{p-1}{np}\le\lambda_k(n,p)\le k\sqrt{p-1}$ for
$k=n+1,n+2,\dots$
\end{enumerate}
Here $\lambda_k(n,p)$ is as in \noderef{LambdaCoef} and $\epsilon_k(n,p)$ is
as in \noderef{EpsCoef}.
\end{lemma}
\begin{proof}
Throughout, $n\ge1$ and $p>1$ are fixed.  Write $\lambda_k=\lambda_k(n,p)$ and
\[
  R_k=4k^2(p-1)+n^2(p-2)^2,
\]
so that, by \noderef{LambdaCoef},
\[
  2\lambda_k=-np+\sqrt{R_k}.
\]
Note $np>0$ because $n\ge1$ and $p>1$, and $p-1>0$.

\emph{Step 0 (the radicand).}  Each of $4k^2(p-1)$ and $n^2(p-2)^2$ is
nonnegative, so $R_k\ge0$ for every $k\ge0$ and $\sqrt{R_k}$ is the ordinary
square root of a nonnegative number.  Moreover, for $k\ge j\ge0$,
\[
  R_k-R_j=4(p-1)(k^2-j^2)\ge0,
\]
with strict inequality when $k>j$, since $p-1>0$.  Hence $k\mapsto R_k$ is
nondecreasing, and $k\mapsto\sqrt{R_k}$ and $k\mapsto\lambda_k$ are
nondecreasing as well, strictly increasing where $R$ strictly increases.

\emph{Step 1 ($\lambda_n=0$).}  Taking $k=n$,
\[
  R_n=4n^2(p-1)+n^2(p-2)^2=n^2\bigl(4p-4+p^2-4p+4\bigr)=n^2p^2 .
\]
Since $np>0$ we get $\sqrt{R_n}=np$, and therefore
$2\lambda_n=-np+np=0$, i.e.\ $\lambda_n=0$.  This is the first half of (23).

\emph{Step 2 ($\epsilon_n=0$).}  By \noderef{EpsCoef} and Step~1,
\[
  \epsilon_n(n,p)=\frac{\lambda_n+n-n}{\lambda_n+n+n}=\frac{0}{2n}=0,
\]
the denominator $2n$ being nonzero because $n\ge1$.  This completes (23).

\emph{Step 3 ($\lambda_k>0$ for $k>n$).}  Let $k\ge n+1$.  By Step~0 and
Step~1,
\[
  R_k-n^2p^2=R_k-R_n=4(p-1)(k^2-n^2)>0,
\]
since $p-1>0$ and $k>n\ge1$ gives $k^2>n^2$.  Hence $R_k>n^2p^2\ge0$, so
$\sqrt{R_k}>\sqrt{n^2p^2}=np$, and therefore $2\lambda_k=-np+\sqrt{R_k}>0$,
i.e.\ $\lambda_k>0$.

\emph{Step 4 (proof of (24)).}  Let $k\ge n+1$ and put
$A=\lambda_k+n$ and $B=k$.  By Step~3, $\lambda_k>0$, and $n\ge1$, so $A>0$;
also $B=k\ge n+1>0$.  Consequently $A+B>0$, and
\[
  A-B<A+B\quad\text{(because $-B<B$, as $B>0$)},\qquad
  -(A+B)<A-B\quad\text{(because $0<2A$)} .
\]
Thus $|A-B|<A+B$, and dividing by $A+B>0$,
\[
  |\epsilon_k(n,p)|=\Bigl|\frac{\lambda_k+n-k}{\lambda_k+n+k}\Bigr|
  =\frac{|A-B|}{A+B}<1 .
\]

\emph{Step 5 (upper bound in (25)).}  Let $k\ge n+1$.  Since $p-1>0$ we have
$\sqrt{p-1}>0$, so $np+2k\sqrt{p-1}>0$; and $\sqrt{R_k}\ge0$.  To conclude
$\sqrt{R_k}\le np+2k\sqrt{p-1}$ it therefore suffices to compare the squares
of these two nonnegative numbers.  Using $(\sqrt{p-1})^2=p-1$,
\[
  \bigl(np+2k\sqrt{p-1}\bigr)^2-R_k
  =n^2p^2+4npk\sqrt{p-1}+4k^2(p-1)-4k^2(p-1)-n^2(p-2)^2
\]
\[
  \qquad=n^2\bigl(p^2-(p-2)^2\bigr)+4npk\sqrt{p-1}
  =4n^2(p-1)+4npk\sqrt{p-1}>0,
\]
where $p^2-(p-2)^2=4(p-1)$ and both displayed terms are positive because
$p-1>0$, $n\ge1$, $k>0$ and $p>0$.  Hence
$\sqrt{R_k}\le np+2k\sqrt{p-1}$, so
\[
  2\lambda_k=-np+\sqrt{R_k}\le 2k\sqrt{p-1},
\]
that is, $\lambda_k\le k\sqrt{p-1}$.

\emph{Step 6 (a lower bound for $\lambda_{n+1}$).}  By Step~1, $R_n=n^2p^2$,
and $\sqrt{R_{n+1}}+np>0$.  Hence
\[
  2\lambda_{n+1}=\sqrt{R_{n+1}}-np
  =\frac{R_{n+1}-n^2p^2}{\sqrt{R_{n+1}}+np}
  =\frac{R_{n+1}-R_n}{\sqrt{R_{n+1}}+np}
  =\frac{4(p-1)(2n+1)}{\sqrt{R_{n+1}}+np},
\]
using $(n+1)^2-n^2=2n+1$ and Step~0.

We next bound the denominator by $2np(2n+1)$, i.e.\ we show
$\sqrt{R_{n+1}}\le np(4n+1)$.  Both sides are nonnegative ($np>0$), so it is
enough to show $R_{n+1}\le n^2p^2(4n+1)^2$.  Now
\[
  n^2p^2(4n+1)^2-R_{n+1}
  =n^2\bigl[\bigl(p(4n+1)\bigr)^2-(p-2)^2\bigr]-4(n+1)^2(p-1)
\]
\[
  \qquad=n^2\bigl[p(4n+1)-(p-2)\bigr]\bigl[p(4n+1)+(p-2)\bigr]-4(n+1)^2(p-1)
\]
\[
  \qquad=n^2(4np+2)\bigl((4n+2)p-2\bigr)-4(n+1)^2(p-1)
  =4\Bigl[n^2(2np+1)\bigl((2n+1)p-1\bigr)-(n+1)^2(p-1)\Bigr],
\]
so it suffices to prove
\begin{equation*}
  n^2(2np+1)\bigl((2n+1)p-1\bigr)\ge(n+1)^2(p-1).
\end{equation*}
Since $p>1$ we have $2np+1>2n+1>0$, and
\[
  (2n+1)p-1=(2n+1)(p-1)+2n\ge(2n+1)(p-1)>0 ,
\]
both factors on the left being positive.  Multiplying the two lower bounds,
\[
  n^2(2np+1)\bigl((2n+1)p-1\bigr)\ge n^2(2n+1)\cdot(2n+1)(p-1)
  =n^2(2n+1)^2(p-1).
\]
Finally $n(2n+1)=2n^2+n\ge n+1$ for $n\ge1$, because $2n^2\ge2>1$; hence
$n^2(2n+1)^2\ge(n+1)^2$, and as $p-1>0$,
\[
  n^2(2n+1)^2(p-1)\ge(n+1)^2(p-1).
\]
This gives the required inequality, hence $\sqrt{R_{n+1}}+np\le 2np(2n+1)$.

Substituting into the displayed formula for $2\lambda_{n+1}$, and using that
its numerator $4(p-1)(2n+1)$ is positive,
\[
  2\lambda_{n+1}=\frac{4(p-1)(2n+1)}{\sqrt{R_{n+1}}+np}
  \ge\frac{4(p-1)(2n+1)}{2np(2n+1)}=\frac{2(p-1)}{np},
\]
that is, $\lambda_{n+1}\ge\dfrac{p-1}{np}$.

\emph{Step 7 (lower bound in (25)).}  Let $k\ge n+1$.  By Step~0,
$\lambda_k\ge\lambda_{n+1}$, and by Step~6,
$\lambda_{n+1}\ge(p-1)/(np)$.  Hence $\lambda_k\ge(p-1)/(np)$.

Steps 1, 2, 4, 5 and 7 are exactly the four assertions (23), (24) and (25).
\end{proof}
